\documentclass[12pt,a4paper]{article}
\usepackage[utf8]{inputenc}
\usepackage[T1]{fontenc}
\usepackage{amsmath,amssymb,amsthm}
\usepackage{physics}
\usepackage{geometry}
\usepackage{enumitem}
\usepackage{hyperref}
\usepackage{fancyhdr}
\usepackage{titlesec}

\geometry{margin=1in}
\pagestyle{fancy}
\fancyhf{}
\rhead{Lecture 3}
\lhead{Advanced Quantum Mechanics}
\cfoot{\thepage}

\titleformat{\section}{\large\bfseries}{}{0em}{}
\titleformat{\subsection}{\normalsize\bfseries}{}{0em}{}

\newtheorem{definition}{Definition}

\title{\textbf{Lecture 3: Central Forces and the Hydrogen Atom} \\ \large Angular Momentum Multiplets, the Coulomb Potential, and the Harmonic Oscillator}
\author{Based on Leonard Susskind's \textit{Advanced Quantum Mechanics}}
\date{}

\begin{document}
\maketitle
\thispagestyle{fancy}

% ============================================================
\section{1. Angular Momentum Multiplets}
% ============================================================

Recall that the simultaneous eigenstates of $L^2$ and $L_z$ are labeled $\ket{\ell, m}$ with:
\[
    L^2\ket{\ell, m} = \hbar^2\ell(\ell+1)\ket{\ell, m}, \qquad L_z\ket{\ell, m} = m\hbar\ket{\ell, m}.
\]
For each value of $\ell$, there are $2\ell + 1$ states with $m = -\ell, -\ell+1, \ldots, \ell$.  These form a \textbf{multiplet}---a complete set of states related by the raising and lowering operators $L_\pm$.

In a rotationally invariant system, the energy depends on $\ell$ but not on $m$, giving a $(2\ell+1)$-fold degeneracy for each angular momentum quantum number.

% ============================================================
\section{2. The Central Force Problem}
% ============================================================

A central force is one in which the potential depends only on the distance from the origin: $V(\vb{r}) = V(r)$.  The Hamiltonian is:
\[
    H = \frac{p^2}{2m} + V(r) = -\frac{\hbar^2}{2m}\nabla^2 + V(r).
\]
Because $V$ depends only on $r$, the system is rotationally invariant: $[H, \vb{L}] = 0$.

\subsection{2.1 Separation of variables}

In spherical coordinates, the Laplacian separates into radial and angular parts:
\[
    \nabla^2 = \frac{1}{r^2}\frac{\partial}{\partial r}\left(r^2\frac{\partial}{\partial r}\right) - \frac{L^2}{\hbar^2 r^2}.
\]
The wave function factorizes as $\psi(r, \theta, \phi) = R(r)\,Y_\ell^m(\theta, \phi)$, where $Y_\ell^m$ are the spherical harmonics---the eigenfunctions of $L^2$ and $L_z$.

\subsection{2.2 The radial equation}

Substituting the separation into the Schr\"odinger equation yields the radial equation:
\[
    -\frac{\hbar^2}{2m}\left[\frac{d^2 u}{dr^2} - \frac{\ell(\ell+1)}{r^2}u\right] + V(r)\,u = E\,u,
\]
where $u(r) = rR(r)$.  The term $\hbar^2\ell(\ell+1)/(2mr^2)$ acts as a \textbf{centrifugal barrier} that pushes particles with higher angular momentum away from the origin.

% ============================================================
\section{3. The Hydrogen Atom}
% ============================================================

For the hydrogen atom, the Coulomb potential is:
\[
    V(r) = -\frac{e^2}{r}.
\]

\subsection{3.1 Energy levels}

The bound-state energies are:
\[
    \boxed{E_n = -\frac{me^4}{2\hbar^2}\frac{1}{n^2} = -\frac{13.6\;\text{eV}}{n^2},}
\]
where $n = 1, 2, 3, \ldots$ is the \textbf{principal quantum number}.

\subsection{3.2 Nodes and quantum numbers}

The radial wave function $R_{n\ell}(r)$ has $n - \ell - 1$ nodes (zeros).  The number of nodes determines how rapidly the wave function oscillates, which is directly connected to the kinetic energy.

For a given $n$, the allowed values of $\ell$ are:
\[
    \ell = 0, 1, 2, \ldots, n-1.
\]

\subsection{3.3 Special degeneracy of the Coulomb potential}

For a general central potential, the energy depends on both $n$ and $\ell$.  The Coulomb potential is special: the energy depends \emph{only} on $n$.  States with different $\ell$ but the same $n$ are degenerate.

The total degeneracy for principal quantum number $n$ is:
\[
    \sum_{\ell=0}^{n-1}(2\ell+1) = n^2.
\]
This ``accidental'' degeneracy arises from a hidden symmetry of the $1/r$ potential (the Laplace--Runge--Lenz vector), beyond ordinary rotational symmetry.

% ============================================================
\section{4. Introduction to the Harmonic Oscillator}
% ============================================================

The harmonic oscillator is one of the most important systems in all of physics.  It appears in:
\begin{itemize}[nosep]
    \item Vibrations of molecules and crystals (phonons).
    \item Modes of the electromagnetic field (photons).
    \item Small oscillations about any stable equilibrium.
    \item The foundation of quantum field theory.
\end{itemize}

\subsection{4.1 Classical harmonic oscillator}

The classical Hamiltonian is:
\[
    H = \frac{p^2}{2m} + \frac{1}{2}m\omega^2 x^2,
\]
where $\omega$ is the angular frequency.  The classical solutions are sinusoidal oscillations at frequency $\omega$.

\subsection{4.2 Quantum treatment preview}

The quantum treatment introduces \textbf{creation} and \textbf{annihilation} operators (also called ladder operators) that algebraically generate the entire energy spectrum:
\[
    E_n = \hbar\omega\left(n + \frac{1}{2}\right), \qquad n = 0, 1, 2, \ldots
\]
The energy levels are equally spaced with separation $\hbar\omega$, and the ground-state energy $E_0 = \tfrac{1}{2}\hbar\omega$ (the \textbf{zero-point energy}) is nonzero.

\medskip
\noindent\textit{The algebraic approach to the harmonic oscillator---using creation and annihilation operators---will be developed fully in the next lecture.  This algebraic machinery will later become the foundation of quantum field theory.}

\end{document}
